%%%%%%%%%%%%%%%%%%%%%%%%%%%%%%%%%%%%%%%%%%%%%%%%%%%%%%%%%%%%%%%%%%%%%%
% Lớp 9 | Chương V. Đường tròn | Mở đầu về đường tròn
% Dạng 9H5?13-2: Tâm đối xứng, trục đối xứng của đường tròn
%%%%%%%%%%%%%%%%%%%%%%%%%%%%%%%%%%%%%%%%%%%%%%%%%%%%%%%%%%%%%%%%%%%%%%


% nguồn: dan_tu_AI (nhập 14/11/2026)
\begin{ex}%[9H5N13-2]%[ma:c0e172b6]
	Đường tròn có bao nhiêu trục đối xứng?
	\choice
	{$1$}
	{$2$}
	{$4$}
	{\True Vô số}
	\loigiai{
		Mỗi đường thẳng đi qua tâm (chứa một đường kính) đều là trục đối xứng của đường tròn, nên đường tròn có vô số trục đối xứng.
	}
\end{ex}

% nguồn: dan_tu_AI (nhập 14/11/2026)
\begin{ex}%[9H5H13-2]%[ma:396c5d55]
	Cho điểm $A$ thuộc đường tròn $(O;4\text{ cm})$. Gọi $A'$ là điểm đối xứng với $A$ qua $O$. Độ dài đoạn thẳng $AA'$ là
	\choice
	{\True $8\text{ cm}$}
	{$4\text{ cm}$}
	{$2\text{ cm}$}
	{$16\text{ cm}$}
	\loigiai{
		$O$ là tâm đối xứng của đường tròn nên $A'$ thuộc đường tròn và $O$ là trung điểm của $AA'$.\\
		Do đó $AA'$ là đường kính, $AA'=2\cdot 4=8\text{ (cm)}$.
	}
\end{ex}

% nguồn: dan_tu_AI (nhập 14/11/2026)
\begin{bt}%[9H5V13-2]%[ma:7c8468fb]
	Cho đường tròn $(O;5\text{ cm})$ và dây $AB=6\text{ cm}$. Gọi $C$, $D$ lần lượt là điểm đối xứng với $A$, $B$ qua $O$. Tính diện tích tứ giác $ABCD$ (đơn vị: cm$^2$).
	\shortans{48}
	\loigiai{
		$O$ là tâm đối xứng của đường tròn nên $C$, $D$ thuộc $(O)$ và $AC$, $BD$ là hai đường kính.\\
		Tứ giác $ABCD$ có hai đường chéo cắt nhau tại trung điểm mỗi đường và bằng nhau ($AC=BD=10$) nên là hình chữ nhật.\\
		Tam giác $ABC$ vuông tại $B$: $BC=\sqrt{AC^2-AB^2}=\sqrt{100-36}=8\text{ (cm)}$.\\
		$S_{ABCD}=AB\cdot BC=6\cdot 8=48\text{ (cm}^2)$.
	}
\end{bt}

% nguồn: dan_tu_AI (nhập 14/11/2026)
\begin{bt}%[9H5V13-2]%[ma:ebf302fd]
	Cho đường tròn $(O;13\text{ cm})$ có đường kính $AB$. Điểm $M$ thuộc đường tròn, $N$ là điểm đối xứng với $M$ qua đường thẳng $AB$ sao cho $MN=24\text{ cm}$. Đoạn $MN$ cắt $AB$ tại $H$, biết $H$ nằm giữa $O$ và $A$. Tính độ dài $AH$ (đơn vị: cm).
	\shortans{8}
	\loigiai{
		$AB$ là trục đối xứng của đường tròn nên $N$ thuộc $(O)$, $AB$ là đường trung trực của $MN$.\\
		Suy ra $MN\perp AB$ tại $H$ và $MH=\dfrac{MN}{2}=12\text{ (cm)}$.\\
		Tam giác $OHM$ vuông tại $H$: $OH=\sqrt{13^2-12^2}=5\text{ (cm)}$.\\
		Vì $H$ nằm giữa $O$ và $A$ nên $AH=OA-OH=13-5=8\text{ (cm)}$.
	}
\end{bt}
